\begin{abstract}
Worldwide image geolocation models such as GeoCLIP estimate GPS coordinates
from a single photograph, yet they provide no account of which visual evidence
supports a prediction. Post-hoc saliency is correlational: it indicates where
activation concentrates without establishing that the highlighted evidence
causes the output. We present GeoProbe, an interactive system for causal
analysis of image geolocation. Given a region specified manually or obtained
from the prompt-free WeDetect-Uni proposal generator, GeoProbe adds
a key-indexed bias to the pre-softmax attention logits of the CLIP vision
tower, independently per layer. It executes GeoCLIP with and without the
intervention and displays both predictions and their geodesic displacement.
On all $2{,}997$ Img2GPS3k images, we search single and multi-layer
interventions using a discovery split and evaluate the selected configuration
on a held-out split. A three-layer configuration achieves $+11.94$~km under
the clipped mean-improvement score on discovery but $-3.66$~km on holdout
(95\% CI $[-20.43,12.40]$), and does not significantly outperform the
selected single layer. Nevertheless, it moves the predicted coordinate by
$189.5$~km on average while retaining an image-embedding cosine similarity of
$0.9984$. The results position GeoProbe as a tool for measuring causal
sensitivity rather than as an accuracy-improvement mechanism, and demonstrate
the importance of held-out validation when interpreting attention
interventions.

\keywords{Image Geolocation \and Explainability \and Attention Intervention
\and Vision--Language Models.}
\end{abstract}
